ગણ $A$ સાન્ત હોય ત્યારે તેનું કદ પ્રાકૃતિક સંખ્યાથી માપી શકાય છે. આ સંખ્યાઓની સરખામણી કરીને કદની સરખામણી થાય છે. અનંત ગણો માટે વાત વધુ જટિલ છે. અનંતનું પ્રથમ સ્તર \emph{ગણનીય અનંત} ગણોનું છે. અખાલી ગણ $A$ના ઘટકોને એક દિશામાં આગળ વધતી અનંત યાદી, એટલે કે \emph{પરિગણના}, તરીકે ગોઠવી શકાય ત્યારે તે ગણનીય છે. બીજા શબ્દોમાં, ત્યારે વ્યાપ્ત ફલન $f\colon\PosInt\to A$ હોય છે. ખાલી ગણ પણ ગણનીય છે. ગણનીય અને અસાન્ત ગણ ગણનીય અનંત છે. કૅન્ટૉરની \emph{વાંકીચૂકી રીત} બતાવે છે કે ગણનીય અનંત ગણોના ઘટકોની જોડોનો ગણ પણ ગણનીય છે. તેથી સંમેય સંખ્યાઓનો ગણ $\Rat$ પણ ગણનીય છે એમ સાબિત થાય છે.

જે અનંત ગણો ગણનીય નથી તે \emph{અગણનીય} છે. ગણ અગણનીય છે એમ સીધું \emph{વિકર્ણ દલીલથી} અથવા \emph{ન્યૂનીકરણથી} બતાવી શકાય છે. વિકર્ણ દલીલમાં $A$ ગણનીય છે એમ ધારીને તેની કાલ્પનિક પરિગણનાથી $A$નો એવો ઘટક વ્યાખ્યાયિત કરીએ છીએ જે બનાવટથી જ પરિગણનાના દરેક ઘટકથી જુદો છે. એટલે એ યાદી આખા $A$ની પરિગણના હોઈ શકે નહીં. ન્યૂનીકરણમાં $A$માંથી કોઈ જાણીતા અગણનીય ગણ $B$ તરફ વ્યાપ્ત ફલન બતાવીએ છીએ. ત્યારે $A$ની કાલ્પનિક પરિગણનાથી $B$ની પરિગણના મળે, જે શક્ય નથી. એટલે $A$ની પરિગણના પણ શક્ય નથી.

સામાન્ય રીતે અનંત ગણોનાં કદ સરખાવી શકાય છે. $A$ અને $B$ વચ્ચે એક-એક વ્યાપ્ત ફલન હોય ત્યારે તેમનું કદ સરખું છે, એટલે કે તેઓ \emph{સમસંખ્યક} છે. $A$માંથી $B$માં એક-એક ફલન હોય તો $A$ એ $B$થી મોટો નથી ($\cardle{A}{B}$). શ્રોડર–બર્નસ્ટાઇન પ્રમેય મુજબ આ સંબંધ અનંત ગણોના કદનો ક્રમ આપે છે. છેવટે, \emph{કૅન્ટૉરનું પ્રમેય} કહે છે કે દરેક $A$ માટે $\cardless{A}{\Pow{A}}$. આ $\Pow{\PosInt}$ અગણનીય છે એ પરિણામનું સામાન્યકરણ છે અને બતાવે છે કે અનંતનાં માત્ર બે સ્તર જ નહીં, અનંત સ્તરો છે.
